% !TEX program = lualatex
\documentclass[12pt, a4paper]{article}
\usepackage{master}

\DeclareWorkGR{Cat}{범주론}{Κατηγορίαι}{Categoriae}
\DeclareWorkGR{Top}{토피카}{Τοπικά}{Topica}
\DeclareWorkGR{DI}{명제론}{Περὶ ἑρμηνείας}{De interpretatione}

\fancyhead[L]{\footnotesize\color{midgray}오르가논 강독}
\fancyhead[R]{\footnotesize\color{midgray}제8주}

\begin{document}
\thispagestyle{firstpage}

\weeksection{제8주}{이사고게 VI--VII장}{부수성과 다섯 술어가능자의 비교}

\subsection{부수성: 오고 가는 것 (VI장)}

\subsubsection{부수성의 정의}

포르피리오스는 부수성(\gr{συμβεβηκός})을 이중적으로 정의합니다. 부수성이란 `주어의 파괴 없이 생겨나고 사라지는 것'이며, `같은 것에 속할 수도 있고 속하지 않을 수도 있는 것'입니다. 두 정의는 서로 다른 측면을 강조합니다 --- 첫 번째는 부수성의 `분리 가능성'을, 두 번째는 부수성의 `우연성'을 강조합니다. 소크라테스가 앉았다가 일어나는 것, 창백해졌다가 다시 혈색이 돌아오는 것 --- 이것들은 모두 소크라테스를 파괴하지 않고 오고 갑니다.

이 정의는 \citework{Cat}{2장}의 `주어 안에 있는 것'(\gr{ἐν ὑποκειμένῳ})과 직접 연결됩니다. 범주론에서 아리스토텔레스는 `주어 안에 있는 것'을 `주어의 부분이 아니면서 주어와 분리되어 존재할 수 없는 것'으로 정의했습니다(1a24--25).

\subsubsection{분리 가능한 부수성과 분리 불가능한 부수성}

포르피리오스는 부수성을 두 가지로 나눕니다.

분리 가능한 부수성(\gr{χωριστὰ συμβεβηκότα})은 잠을 잠, 움직임, 앉아 있음처럼 어렵지 않게 주어로부터 분리되는 것입니다. 분리 불가능한 부수성(\gr{ἀχώριστα συμβεβηκότα})은 까마귀의 검음, 에티오피아인의 검음처럼 사실상 항상 주어에 속하지만, 관념적으로는(\gr{κατ᾽ ἐπίνοιαν}) 분리 가능한 것입니다. 포르피리오스의 검증 기준: `우리는 흰 까마귀를 생각할 수 있다 --- 주어의 파괴 없이(\gr{ἄνευ τῆς τοῦ ὑποκειμένου φθορᾶς}).' 이 기준은 자연적 필연성이 아니라 논리적 가능성입니다 --- 물리 세계에서 실제로 발생하느냐가 아니라, 모순 없이 생각할 수 있느냐입니다.

\subsubsection{종차, 고유성, 분리 불가능한 부수성의 삼자 비교}

세 가지 모두 주어에 항상(또는 사실상 항상) 속하지만, 논리적 지위가 다릅니다.

{\footnotesize
\begin{concepttable}{|c|c|c|l|}
\hline
\rowcolor{tableheadblue}
{\color{h1navy}\bfseries 항목} &
{\color{h1navy}\bfseries 본질적?} &
{\color{h1navy}\bfseries 역전?} &
{\color{h1navy}\bfseries 설명} \\ \hline
종차 (이성적인) & 예 & 아니오 & 류 + 종차 = 정의 \\ \hline
고유성 (웃을 수 있음) & 아니오 & 예 & 본질에서 따라나오나 본질 아님 \\ \hline
분리불가 부수성 (검음) & 아니오 & 아니오 & 항상 속하나 논리적 분리 가능 \\ \hline
\end{concepttable}
}

종차는 본질적입니다 --- `이성적인'을 제거하면 사람은 더 이상 사람이 아닙니다. 고유성은 비본질적이지만 역전됩니다 --- 웃을 수 있음은 사람에게만, 모든 사람에게, 항상 속합니다. 분리 불가능한 부수성은 비본질적이고 역전되지도 않습니다 --- 검음은 까마귀에 항상 속하지만 까마귀에게만 속하는 것이 아닙니다. 추가로, 본질적 종차는 정도의 차이를 허용하지 않지만, 부수성은 정도의 차이를 허용할 수 있습니다 --- 이 까마귀가 저 까마귀보다 `더 검을' 수 있습니다.

\subsubsection{범주론의 사중 분류와의 대응}

이사고게의 다섯 술어가능자 체계가 범주론의 존재론적 틀과 어떻게 대응하는지를 정리하면, 두 텍스트의 관계가 분명해집니다. 범주론의 `주어에 대해 말해지는 것'(said-of)은 본질적 술어화에 해당하며, 이사고게의 류, 종차, 종이 이에 속합니다. 범주론의 `주어 안에 있는 것'(present-in)은 비본질적 술어화에 해당하며, 부수성이 이에 속합니다. 고유성은 흥미로운 중간 지대를 차지합니다 --- 비본질적이지만 필연적입니다. 범주론의 사중 분류에서 고유성은 `주어에 대해 말해지면서 주어 안에 있는 것'(보편적 속성)에 가장 가깝습니다.

\subsection{다섯 술어가능자의 비교 (VII장)}

\subsubsection{공통점과 핵심 쌍별 비교}

모든 술어가능자가 공유하는 가장 근본적인 공통점은 여럿에 대해 술어된다(\gr{κατὰ πλειόνων κατηγορεῖται})는 것입니다 --- 이것이 보편자(\gr{τὰ καθόλου})의 정의적 특징입니다.

류와 종차의 비교: 둘 다 종에 대해 본질적으로 술어되지만, 류는 `무엇인지'(\gr{τί ἐστι})를 말해주고, 종차는 `어떤 종류의 것인지'(\gr{ὁποῖόν τί ἐστι})를 말해줍니다. 류와 고유성의 비교: 류는 종 전체에 술어되고 역전되지 않지만, 고유성은 종과 역전됩니다. 종차와 고유성의 비교: 둘 다 한 종에 속하지만, 종차는 본질적이고 고유성은 비본질적입니다.

\subsubsection{다섯 술어가능자의 전체 구조}

{\footnotesize
\begin{concepttable}{|c|c|c|c|c|}
\hline
\rowcolor{tableheadblue}
{\color{h1navy}\bfseries 술어가능자} &
{\color{h1navy}\bfseries 본질적?} &
{\color{h1navy}\bfseries 역전?} &
{\color{h1navy}\bfseries \gr{τί ἐστι}?} &
{\color{h1navy}\bfseries 외연(종 대비)} \\ \hline
류 & 예 & 아니오 & 예 & 넓음 \\ \hline
종 & 예 & — & 예 & (기준) \\ \hline
종차 & 예 & 아니오 & 아니오* & 넓음 \\ \hline
고유성 & 아니오 & 예 & 아니오 & 동일 \\ \hline
부수성 & 아니오 & 아니오 & 아니오 & 다양 \\ \hline
\end{concepttable}
}

{\footnotesize * 종차는 \gr{ὁποῖόν τί ἐστι} --- `어떤 종류의 것인지'를 말해줌.\par}

\subsection{이사고게의 역사적 의의}

\subsubsection{구논리학의 첫 번째 텍스트}

이사고게는 `구논리학'(logica vetus)의 첫 번째 구성 요소가 되었습니다. 1주차에서 다루었듯이, 6세기부터 12세기 중반까지 라틴 서방에서 접근 가능했던 논리학 텍스트는 이사고게, 범주론, 명제론(모두 보에티우스의 번역), 그리고 보에티우스 자신의 논리학 논고들뿐이었습니다. 보에티우스 자신이 `포르피리오스 이후에 논리학에 착수한 모든 사람이 이 책으로 시작했다'고 기록하고 있습니다.\footnote{Boethius, \gri{In Isagogen Porphyrii commenta}, editio prima, 12.20--21.}

\subsubsection{보에티우스의 이중 주석과 보편자 논쟁}

보에티우스(c.\ 480--524)는 이사고게에 대해 두 편의 주석서를 썼습니다. 제1주석서(c.\ 505)는 마리우스 빅토리누스의 번역에 기초한 초급용 대화편이고, 제2주석서(c.\ 508)는 보에티우스 자신의 번역에 기초한 5권짜리 고급용 저술입니다.\footnote{S.\ Brandt, ed., \gri{Anicii Manlii Severini Boetii in Isagogen Porphyrii Commenta}, CSEL 38 (Vienna, 1906).} 제2주석서의 제1권은 포르피리오스 텍스트 불과 한 쪽 남짓에 대해 전체 분량을 할애합니다.

보에티우스의 해법은 추상론(abstractionism)입니다: 마음이 실재에서 분리 불가능한 것을 사유에서 분리하여 류와 종의 개념을 형성한다는 것입니다. 그러나 그는 플라톤적 입장도 명시적으로 열어 둡니다: `나는 그들의 견해 사이에서 결정을 내리는 것이 적절하지 않다고 판단했다.' 이 외교적 유보 --- 포르피리오스의 거부에 보에티우스의 거부가 겹쳐진 것 --- 가 중세 보편자 논쟁의 구조적 틀을 만들었습니다.\footnote{Paul Vincent Spade, \gri{Five Texts on the Mediaeval Problem of Universals} (Indianapolis: Hackett, 1994), 20--25.}

\subsubsection{보편자 논쟁의 전개}

로스켈리누스(Roscelin, c.\ 1050--c.\ 1125)는 보편자가 `목소리의 내뿜음'(flatus vocis)에 불과하다고 주장했다고 전해집니다 --- 극단적 유명론입니다. 아벨라르(Abelard, 1079--1142)는 양쪽 극단을 모두 비판하고, 보편자는 의미 있는 단어라고 주장했습니다. 둔스 스코투스(Duns Scotus, c.\ 1266--1308)는 `공통 본성'(natura communis)과 `이것임'(haecceitas)의 결합이라는 정교한 실재론을, 오캄(William of Ockham, c.\ 1287--1347)은 보편자가 마음의 개념에 불과하다는 유명론을 발전시켰습니다.\footnote{SEP, ``The Medieval Problem of Universals'' (Gyula Klima, rev.\ 2022).}

\subsubsection{이슬람 세계의 전승과 이사고게의 유산}

아랍어화된 이름 이사구지(\textit{Is\=agh\=uj\={\i}})로 이 텍스트는 이슬람 세계 전역에서 표준 논리학 입문서로 남았습니다. 아씨르 알딘 알아브하리(Ath\={\i}r al-D\={\i}n al-Abhar\={\i})의 독립 저술 이사구지는 20세기까지 북아프리카와 중동의 교육 기관에서 사용되었습니다.

이사고게의 의의는 세 가지로 요약할 수 있습니다. 첫째, 생성적 모호성 --- 자신의 세 물음에 대답하지 않음으로써 다양한 전통이 수용할 수 있는 틀을 만들었습니다. 둘째, 체계적 우아함 --- 다섯 술어가능자 체계는 교육 전통에서 아리스토텔레스의 네 가지 분류를 대체했습니다. 셋째, 제도적 내구성 --- 세 문명권에 걸쳐 천 년 이상 철학 교육의 관문으로 자리 잡으면서, 분류, 본질, 술어화에 대한 사고 습관을 형성했습니다.

다음 주(9주차)부터 우리는 명제론으로 넘어갑니다. 범주론이 비결합 표현들을 분석했다면, 명제론은 결합 표현들(명제)을 분석합니다 --- 참이거나 거짓이 될 수 있는 진술의 논리적 구조를 탐구하는 것입니다. 3주차에서 다루었던 `진술이 참에서 거짓으로 바뀔 때 변하는 것은 무엇인가'라는 물음, 그리고 5주차의 `긍정과 부정의 대립'이라는 개념이 명제론의 출발점이 됩니다.

% ── 참고문헌 ──
\subsection*{참고문헌}
\begin{references}

\refitem Barnes, Jonathan. \gri{Porphyry: Introduction}. Clarendon Later Ancient Philosophers. Oxford: Clarendon Press, 2003.

\refitem Brandt, Samuel, ed. \gri{Anicii Manlii Severini Boetii in Isagogen Porphyrii Commenta}. CSEL 38. Vienna, 1906.

\refitem Busse, Adolf, ed. \gri{Porphyrii Isagoge et in Aristotelis Categorias Commentarium}. CAG IV.1. Berlin: G.\ Reimer, 1887. Eng.\ trans.: Jonathan Barnes, \gri{Porphyry: Introduction}, Clarendon Later Ancient Philosophers. Oxford: Clarendon Press, 2003.

\refitem Emilsson, Eyjólfur. ``Porphyry.'' In \gri{Stanford Encyclopedia of Philosophy}, rev.\ 2021.\newline https://plato.stanford.edu/entries/porphyry/.

\refitem Klima, Gyula. ``The Medieval Problem of Universals.'' In \gri{Stanford Encyclopedia of Philosophy}, rev.\ 2022.\newline https://plato.stanford.edu/entries/universals-medieval/.

\refitem Spade, Paul Vincent. \gri{Five Texts on the Mediaeval Problem of Universals}. Indianapolis: Hackett, 1994.

\end{references}

\end{document}
